\documentclass{sajzk}

\usepackage{array}

\usepackage{minted}
% \usepackage{listings}
% \lstset{basicstyle=\ttfamily}
% \usepackage{xcolor}
% \usepackage{mdframed}
% 
% \definecolor{light-gray}{gray}{0.95} % the shade of grey that stack exchange uses

\usetikzlibrary {arrows.meta, angles, quotes}
\usetikzlibrary {calc}

\begin{document}
\section{Gesetz vom ausgeschlossenen Widerspruch (GAW)}
\label{3z0c}
\lhead{:mathe:logik:kalkül:}
\textbf{GAW:} $\vdash \lnot(A\land\lnot A)$

\textbf{Theorem:} $\vdash \lnot(P\land\lnot P)$

\textbf{Ableitung:}
\begin{center}
\begin{tabular}{|c|c|c|c|}
  \hline
  Annahme            & Nr. & Formel                    & Regel \\
  \hline
  $1^\star$          & (1)    & $P\land\lnot P$        & AE \\
  \hline
                     & (2)    & $\lnot(P\land\lnot P)$ & 1,1 RAA \\
  \hline
\end{tabular}
\end{center}

\subsection{Referenzen}
\begin{itemize}
    \item \href{s52i.pdf}{Logische Gesetze} s52i.tex
\end{itemize}

\subsection{Literatur}
\begin{itemize}
    \item Timm Lampert - Klassische Logik (2004) Seite 121
\end{itemize}
\end{document}
